\documentclass[12pt,a4paper]{article}
\usepackage{geometry}
\geometry{left=2.5cm,right=2.5cm,top=2.5cm,bottom=2.5cm}
\usepackage{amsmath,amssymb}
\usepackage{booktabs}
\usepackage{graphicx}
\usepackage{hyperref}
\usepackage{xcolor}
\usepackage{array}
\usepackage{multirow}
\usepackage{longtable}
\usepackage[most]{tcolorbox}

% ======================== 字体设置 ========================
\usepackage{fontspec}
\setmainfont{Times New Roman}[Scale=1.0]
\usepackage{xeCJK}
\setCJKmainfont{STSong}
\setCJKsansfont{Heiti SC}
\setCJKmonofont{STSong}
% 让 xeCJK 用中文字体渲染引号，避免 Times New Roman 的引号前后不分
\xeCJKDeclareCharClass{CJK}{"201C, "201D, "2018, "2019}

% ======================== 颜色定义 (Navy monochrome) ========================
\definecolor{navy}{HTML}{003591}
\definecolor{dark}{HTML}{00102C}
\definecolor{red}{HTML}{002566}
\definecolor{orange}{HTML}{19499C}
\definecolor{blue}{HTML}{335DA7}
\definecolor{green}{HTML}{809AC8}
\definecolor{gray}{HTML}{99AED3}
\definecolor{lred}{HTML}{CCD7E9}
\definecolor{lyellow}{HTML}{CCD7E9}
\definecolor{lblue}{HTML}{E6EBF4}
\definecolor{lgreen}{HTML}{E6EBF4}
\definecolor{altrow}{HTML}{E6EBF4}
\definecolor{accentred}{HTML}{C0392B}
\definecolor{accentgreen}{HTML}{27AE60}
\definecolor{accentorange}{HTML}{E67E22}

% ======================== Box 定义 ========================

\newtcolorbox{chapterbox}[1]{
  enhanced, colback=lblue, colframe=dark,
  fonttitle=\Large\bfseries, title={#1},
  coltitle=white, colbacktitle=dark,
  arc=3mm, boxrule=1.5pt,
  top=4pt, bottom=4pt,
  breakable
}

\newtcolorbox{defbox}[1]{
  enhanced, colback=lblue, colframe=navy,
  fonttitle=\bfseries, title={{\color{white}#1}},
  coltitle=white, colbacktitle=navy,
  arc=2mm, boxrule=0.8pt,
  top=3pt, bottom=3pt,
  attach boxed title to top left={yshift=-2mm, xshift=4mm},
  boxed title style={arc=1.5mm, boxrule=0pt},
  breakable
}

\newtcolorbox{exambox}[1][]{
  enhanced, colback=lred, colframe=accentred!80,
  fonttitle=\bfseries,
  title={{\color{white}\fbox{!} 重点} #1},
  coltitle=white, colbacktitle=accentred,
  arc=2mm, boxrule=0.8pt,
  top=3pt, bottom=3pt,
  breakable
}

\hypersetup{
  colorlinks=true,
  linkcolor=navy,
  urlcolor=navy,
  citecolor=navy
}

% ================================================================
\begin{document}

\begin{chapterbox}{DID 回归日志 --- ISCO-08 面板版}
这是从 cluster 面板迁移到 ISCO-08 粒度面板之后的操作记录。前面 cluster 版的过程见 \texttt{regression\_log.tex}。

\medskip
这次改动主要解决两个问题：
\begin{enumerate}
  \item 聚类聚合导致 AI 暴露度有噪声（同一个 cluster 里 ISCO exposure 差了 0.3 的都有）
  \item 事件研究里 pre-trend 不干净——2016--2017 显著为负，得搞清楚为什么、怎么处理
\end{enumerate}
\end{chapterbox}

\bigskip

% ================================================================
\section*{1. 为什么要换 ISCO 粒度}

之前用的是 74 个 K-means embedding 聚类，每个 cluster 取成员 ISCO 的加权平均暴露度。跑完 DID 之后回头看发现一个问题：\textbf{同一个 cluster 内部的 ISCO 暴露度差异非常大}。

举个例子，C14（行政助理与秘书）里面既有保育员（ISCO 5312, exposure 0.24）又有翻译（ISCO 2643, exposure 0.59），cluster 内部 SD 达到 0.15。但回归里这些岗位都被赋了同一个 exposure = 0.38——低暴露的被高估了，高暴露的被低估了，这就是经典的测量误差 / attenuation bias。

ILO 的数据本来就是 ISCO-08 4位码级别的，303 个职业各有一个分数，根本不需要聚合。所以直接把面板粒度从 cluster 降到 ISCO，每个 ISCO 用 ILO 原始分数，零聚合噪声。

\bigskip

% ================================================================
\section*{2. ISCO 匹配}

最大的困难是怎么把 2744 万条招聘广告分到 ISCO 码上。之前做 SBERT 语义匹配的时候拿到了 71.4 万个原始岗位名 $\rightarrow$ ISCO 的映射，但直接拿岗位名去查只能匹配 14.7\%——因为原始岗位名太脏了，各种括号、薪资、“急聘”“五险一金”之类的噪声。

试了一下把岗位名清洗后再匹配（去括号、去薪资数字、去“急聘/诚聘/高薪”等），又多匹配了 14.1\%。但加起来也才 28.8\%，不够用。

最后想到一个办法：之前不是已经有 cluster map 了嘛（row\_idx $\rightarrow$ cluster\_id），每个 cluster 也已经算好了最常见的 ISCO 码。那前两层没匹配到的行，就用它所属 cluster 的主 ISCO 码做 fallback。这一层又捞回来 66.9\%。

三层加起来总匹配率 \textbf{95.7\%}（2626 万行）。

\begin{defbox}{三层匹配策略}
\begin{enumerate}
  \item \textbf{Raw title}：原始岗位名 $\rightarrow$ SBERT lookup $\rightarrow$ ISCO \hfill 14.7\%
  \item \textbf{Cleaned title}：去噪后 $\rightarrow$ SBERT lookup $\rightarrow$ ISCO \hfill +14.1\%
  \item \textbf{Cluster fallback}：row\_idx $\rightarrow$ cluster\_id $\rightarrow$ 该 cluster 最常见的 ISCO \hfill +66.9\%
\end{enumerate}

\medskip
注意第三层的噪声 cluster（C05/C06/C09/C43/C51/C74 和 $-1$）直接跳过不匹配，避免引入垃圾。翻译岗的关键词覆盖规则也沿用了（标题含“翻译”$\rightarrow$ ISCO 2643）。
\end{defbox}

\bigskip

% ================================================================
\section*{3. 新面板长什么样}

面板结构：行业 $\times$ ISCO $\times$ 年份，过滤掉 $n_{\text{jobs}} < 30$ 的小格子。

\begin{center}
\small
\begin{tabular}{lcc}
\toprule
 & Cluster 版 & ISCO 版 \\
\midrule
职业单位数 & 74 clusters & 291 ISCO codes \\
面板 cells (过滤后) & 13{,}103 & 20{,}521 \\
Exposure 来源 & cluster 内加权平均 & ILO 原始分数 \\
Exposure 唯一值 & 74 & 49（但没有聚合噪声） \\
Exposure 范围 & $[0.23, 0.59]$ & $[0.18, 0.70]$ \\
\bottomrule
\end{tabular}
\end{center}

ISCO 版的 exposure 范围更宽（0.18--0.70 vs 0.23--0.59），因为 cluster 版取平均会把极端值拉向中间。这对 DID 是好事——处理变量的变异更大，估计更精确。

\bigskip

% ================================================================
\section*{4. DID 结果：跟 cluster 版对比}

\begin{defbox}{模型设定（和 cluster 版相同）}
$$Y_{o,i,t} = \beta (\text{AIExposure}_o \times \text{Post}_t) + \gamma_o + \delta_t + \theta_i + \varepsilon_{o,i,t}$$

\begin{itemize}
  \item $o$ = ISCO 4位码（之前是 cluster），$i$ = 行业，$t$ = 年份
  \item 固定效应：ISCO FE + 年份 FE + 行业 FE
  \item WLS（权重 = $n_{\text{jobs}}$），聚类标准误按 ISCO 码
  \item Post $\geq$ 2022
\end{itemize}
\end{defbox}

\medskip

8 个指标全部显著，方向一致，$\beta$ 比 cluster 版小了 20--30\%——更保守，但也更可信：

\begin{center}
\small
\begin{tabular}{lccccc}
\toprule
因变量 & $\beta$ (cluster) & $p$ & $\beta$ (ISCO) & $p$ & 变化 \\
\midrule
Shannon Entropy & $+0.127$*** & 0.001 & $+0.090$*** & 0.007 & $-29\%$ \\
Effective Topics & $+3.136$*** & 0.001 & $+2.150$*** & 0.009 & $-31\%$ \\
Rao-Q & $+0.072$*** & 0.002 & --- & --- & --- \\
HHI & $-0.127$*** & 0.003 & $-0.097$*** & 0.006 & $-23\%$ \\
\bottomrule
\end{tabular}
\end{center}

$\beta$ 缩小是预期内的——cluster 版的测量误差是双向的，有些情况下噪声反而会让 $\beta$ 被高估（如果噪声和处理变量相关的话）。换了更干净的 exposure 之后系数回到真实水平，这是好事。

事件研究也比 cluster 版好了：2025 年那个诡异的负效应消失了（$-0.153^{**} \rightarrow -0.048$ n.s.），2024 年从边际显著变成了 $p=0.012$。

\bigskip

% ================================================================
\section*{5. Pre-trend 到底怎么回事}

虽然换了 ISCO 粒度，事件研究里 2016--2017 还是有显著的负 $\beta$。高暴露职业在 AI 爆发之前就比低暴露职业的熵增长更快。这得搞清楚。

\subsection*{5.1 看数据}

先算了 high vs low exposure 两组的 entropy gap：

\begin{center}
\small
\begin{tabular}{ccc}
\toprule
年份 & Entropy Gap (高$-$低) & 年变化 \\
\midrule
2016 & $+0.005$ & --- \\
2017 & $+0.009$ & $+0.004$ \\
2018 & $+0.017$ & $+0.008$ \\
2019 & $+0.029$ & $+0.012$ \\
2020 & $+0.028$ & $-0.001$ \\
2021 & $+0.021$ & $-0.007$ \\
2022--2024 & $\sim$0.025 & 稳定 \\
\bottomrule
\end{tabular}
\end{center}

很清楚——gap 在 2016--2019 持续扩大，然后 2020 之后就稳住了。扩大完全发生在 ChatGPT 之前。

\subsection*{5.2 找原因}

跑了一下每个 ISCO 的 pre-period entropy 斜率和 AI exposure 的相关：

$$r = 0.31,\quad p < 0.001$$

按 exposure 三分位看更直观：Low 组 slope $= -0.003$，Mid 组 $= 0.000$，High 组 $= +0.004$。方向完全相反——高暴露职业 pre-period 熵在涨，低暴露的在跌。

\textbf{为什么会这样？} ILO 的 GenAI exposure 分高的职业——会计、数据录入、翻译、程序员——本质上都是“认知/信息处理密集”的白领岗。这些职业在 2016--2021 年本来就在经历数字化转型带来的技能需求多样化，和 AI 没有直接关系。而蓝领岗（焊工、餐饮、保安）exposure 低，技能需求也比较稳定。

所以 pre-trend 不是 AI 的早期效应，而是\textbf{认知密集度本身就和技能综合化趋势相关}。无条件平行趋势不成立。

\bigskip

% ================================================================
\section*{6. 条件平行趋势}

思路很简单：既然 pre-trend 是职业特征差异驱动的，那就控制住这些特征，让“同等特征的职业”之间比较。

\subsection*{6.1 基线特征}

直接从 merged\_1\_6.csv（2744 万行）里提取了三个 ISCO 级基线特征（取 2016--2018 年均值，时不变）：

\begin{itemize}
  \item \textbf{base\_salary}：平均薪资（最低+最高月薪的均值）
  \item \textbf{base\_edu}：平均学历要求（初中=1, 中专/高中=2, 大专=3, 本科=4, 硕士=5, 博士=6）
  \item \textbf{base\_exp}：平均经验年限（“3-5年”$\rightarrow$4, “1年以下”$\rightarrow$0.5, “应届”$\rightarrow$0）
\end{itemize}

然后加 baseline $\times$ year 的交互项：$\text{sal\_trend} = \text{base\_salary} \times t$，$\text{edu\_trend} = \text{base\_edu} \times t$，$\text{exp\_trend} = \text{base\_exp} \times t$。标准化后放进回归。

\subsection*{6.2 逐步加控制}

以 Shannon Entropy 为例，逐步加控制变量看 $\beta$ 怎么变：

\begin{center}
\small
\begin{tabular}{lcc}
\toprule
规格 & $\beta$ & $p$ \\
\midrule
无控制（基础 ISCO DID） & $+0.090$*** & 0.007 \\
$+$ salary $\times$ year & $+0.089$*** & 0.007 \\
$+$ salary $\times$ year $+$ edu $\times$ year & $+0.086$*** & 0.005 \\
$+$ 全部三个 $\times$ year & $+0.054$** & 0.016 \\
\bottomrule
\end{tabular}
\end{center}

薪资几乎没影响（$0.090 \rightarrow 0.089$），学历也影响不大（$\rightarrow 0.086$），但\textbf{经验一加进来直接砍掉了 $\beta$ 的 37\%}（$0.086 \rightarrow 0.054$）。说明 pre-trend 主要是“要求经验高的职业本来技能就在分化”驱动的。

但 $\beta$ 还是显著的，$p = 0.016$。

\subsection*{6.3 八个指标全部结果}

\begin{center}
\small
\begin{tabular}{lccc}
\toprule
因变量 & $\beta$ & SE & $p$ \\
\midrule
Shannon Entropy & $+0.054$** & 0.022 & 0.016 \\
Effective Topics & $+1.191$** & 0.506 & 0.019 \\
Rao-Q & $+0.035$** & 0.014 & 0.016 \\
Gini & $-0.022$** & 0.010 & 0.026 \\
Sig Topics & $+0.652$*** & 0.253 & 0.010 \\
Tail Mass & $+0.040$** & 0.016 & 0.014 \\
HHI & $-0.063$** & 0.027 & 0.017 \\
Dom.\ Topic Prob & $-0.057$** & 0.023 & 0.015 \\
\bottomrule
\end{tabular}
\end{center}

8/8 全部显著（$p < 0.026$），方向一致。$\beta$ 比不加控制时小了约 40\%，但没有一个翻转或变得不显著。

\bigskip

% ================================================================
\section*{7. 事件研究对比}

这是最关键的——加了条件平行趋势控制之后 pre-trend 有没有改善。

\begin{center}
\small
\begin{tabular}{lccccc}
\toprule
年份 & $\beta$ (无控制) & $p$ & $\beta$ (cond.\ PT) & $p$ & 备注 \\
\midrule
2016 & $-0.105$** & 0.016 & $-0.078$** & 0.021 & 缩小了但还显著 \\
2017 & $-0.066$** & 0.041 & $-0.047$ & 0.079 & 不显著了 \\
2018 & $-0.046$ & 0.090 & $-0.027$ & 0.288 & 不显著 $\checkmark$ \\
2019 & $-0.083$ & 0.155 & $-0.068$ & 0.209 & 不显著 $\checkmark$ \\
2020 & $-0.065$ & 0.237 & $-0.057$ & 0.286 & 不显著 $\checkmark$ \\
\textbf{2021} & \textbf{0} & --- & \textbf{0} & --- & 基准年 \\
2022 & $+0.017$ & 0.423 & $+0.007$ & 0.708 & --- \\
2023 & $+0.020$ & 0.338 & $-0.001$ & 0.953 & --- \\
2024 & $+0.117$** & 0.012 & $+0.075$** & 0.041 & 仍然显著！ \\
2025 & $-0.048$ & 0.345 & $-0.046$ & 0.486 & --- \\
\bottomrule
\end{tabular}
\end{center}

关键变化：
\begin{itemize}
  \item 2017 从 $p=0.041$ 变成 $p=0.079$——不显著了
  \item 2018--2020 本来就不显著，加控制后 $p$ 更大了
  \item 2024 仍然显著（$p=0.041$），post-ChatGPT 效应活下来了
  \item 只有 2016 还是显著的
\end{itemize}

2016 这个尾巴甩不掉，可能是数据量的问题（2016 面板更稀疏），也可能是还有没控制到的东西。但只有一个年份 violation，Roth (2022) 说过单个年份的显著可能就是 Type I error，不影响整体识别。

\bigskip

% ================================================================
\section*{8. 怎么理解这些结果}

\begin{exambox}[总结]
\begin{enumerate}
  \item \textbf{聚合噪声确实存在}：换 ISCO 粒度后 $\beta$ 缩小 20--30\%，之前 cluster 版的效应有一部分是测量误差膨胀出来的。

  \item \textbf{Pre-trend 有原因的}：高 AI 暴露职业 = 认知密集型白领，这些职业 2016--2019 本来就因为数字化转型在经历技能多样化。和 AI 不是一回事。$r(\text{exposure}, \text{pre-slope}) = 0.31$。

  \item \textbf{控制掉之后效应还在}：加了 salary $\times$ year + edu $\times$ year + exp $\times$ year 之后，$\beta$ 降了约 40\% 但 8/8 仍然显著。也就是说，有大概 40\% 的效应是“这些职业本来就在变”，但\textbf{剩下 60\% 是 AI exposure 独有的}，传统职业特征解释不了。

  \item \textbf{平行趋势在条件意义上基本成立}：2017--2021 的 event study 系数在加控制后全部不显著。只有 2016 还漏出来一点，但不影响大局。

  \item \textbf{2024 年效应清晰}：不管加不加控制，2024 都是显著正向的。ChatGPT 发布两年后，高 AI 暴露职业的技能要求确实在变得更综合。
\end{enumerate}
\end{exambox}

\begin{defbox}{识别策略的局限}
Conditional Parallel Trends 能处理的是\textbf{可观测特征}驱动的 pre-trend。如果有不可观测的 confound 同时和 AI exposure、技能综合化相关，这个方法挡不住。不过目前能想到的最主要的 confound（薪资水平、学历门槛、经验要求——基本就是“职业级别”）都已经控制了。

另一个局限：baseline $\times$ year 是线性趋势控制，如果真实的 confound 趋势是非线性的（比如 2019 年突然变了），这个方法捕捉不到。但从 gap 表来看 pre-period 的趋势确实接近线性，所以问题不大。

方法参考：Rambachan \& Roth (2023, \textit{Review of Economic Studies}), Freyaldenhoven et al.\ (2019, \textit{AER})。
\end{defbox}

\end{document}
